\chapter{Cancelamento do registro profissional}
\label{chap:pf-cancelamento}

\section{O que é?}

Cancelamento de registro de psicóloga/o. O registro cancelado impede o exercício profissional.

\section{Quem pode utilizar esse serviço?}

Psicólogas/os inscritas/os e ativas/os, que não exercerão a profissão pelo período em que o registro estiver cancelado.

\section{Como solicitar o serviço?}

A solicitação de cancelamento de inscrição deverá ser enviada pelo correio com aviso de recebimento (AR), acompanhada de carteira de identidade profissional original e do formulário de cancelamento,disponível no \href{https://www.crpsp.org/pagina/view/312}{\emph{site}do CRP~SP}, devidamente assinado.

\section{Que informações ou documentos são necessários?}

\begin{itemize}
  \item  Requerimento  de cancelamento de inscrição (pessoa física) preenchido, datado e  assinado;
  \item  CIP original.
\end{itemize}

\section{Quais são as etapas para a realização desse serviço?}

\begin{itemize}
  \item  Envio da documentação necessária via carta com aviso de recebimento. O  registro será considerado cancelado a partir da data de registro da  correspondência no correio; ou
  \item  entrega dos documentos necessários na subsede de referência (confira  os endereços \href{https://www.crpsp.org/sede/index}{no \emph{site}})
\end{itemize}

\section{Quanto custa?}

Sem custo.

\noindent Obs.: O cancelamento não desobriga a/o psicóloga/o do pagamento de débitos anteriores à data de cancelamento.

\section{Quanto tempo leva?}

30 dias para os processos enviados via correio; imediato no caso de solicitação presencial.

\section{Como ter acesso ao atual \emph{status} do serviço?}

Pode-se consultar a inscrição no \href{https://cadastro.cfp.org.br/}{Cadastro Nacional de Profissionais de Psicologia}.

\section{Legislação relacionada}

\begin{itemize}
  \item  \href{https://atosoficiais.com.br/cfp/resolucao-administrativa-financeira-n-3-2007-institui-a-consolidacao-das-resolucoes-do-conselho-federal-de-psicologia?origin=instituicao&q=resolu\%C3\%A7\%C3\%A3o\%203/2007}{Resolução  CFP nº~03/2007} (Consolidação das Resoluções do Conselho Federal de  Psicologia);
  \item  \href{https://transparencia.cfp.org.br/wp-content/uploads/sites/7/2020/05/Manual-de-Procedimentos-Administrativos-e-Financeiro-CFP-RES-20.2018.pdf}{\href{https://atosoficiais.com.br/cfp/resolucao-administrativa-financeira-n-20-2018-altera-o-manual-de-procedimentos-administrativo-e-financeiro-do-sistema-conselhos-de-psicologia-anexo-da-resolucao-cfp-n-202018-a-resolucao-cfp-n-03-2007-a-resolucao-cfp-n-16-2019-e-da-outras-providencias}{Resolução  CFP nº~20/2018}}, que determina a revisão e ampliação do  \href{https://transparencia.cfp.org.br/wp-content/uploads/sites/7/2020/05/Manual-de-Procedimentos-Administrativos-e-Financeiro-CFP-RES-20.2018.pdf}{\emph{Manual  de procedimentos administrativos e financeiros}}
\end{itemize}

\sumario
